% Generated draft: requires mathematical and source review.
\documentclass[11pt]{article}
\usepackage[margin=25mm]{geometry}
\usepackage{fontspec}
\setmainfont{Latin Modern Roman}
\usepackage{amsmath,amssymb,mathtools,unicode-math}
\setmathfont{Latin Modern Math}
\usepackage{enumitem,needspace,xurl,hyperref}
\setlength{\parindent}{0pt}
\setlength{\parskip}{5pt}
\setlength{\emergencystretch}{3em}
\newcommand{\N}{\mathbb{N}}
\newcommand{\Z}{\mathbb{Z}}
\newcommand{\Q}{\mathbb{Q}}
\newcommand{\R}{\mathbb{R}}
\newcommand{\C}{\mathbb{C}}
\newcommand{\F}{\mathbb{F}}
\newcommand{\A}{\mathbb{A}}
\newcommand{\PP}{\mathbb{P}}
\newcommand{\E}{\mathbb{E}}
\newcommand{\eps}{\varepsilon}
\newcommand{\dd}{\,\mathrm{d}}
\begin{document}
\section*{UnsolvedMath ID: 3900006}
\textbf{Source problem code:} catalogue entry
\subsection{Definitions and mathematical statement}
Use the upper half-plane model
\[
\mathbb H^2=\{(x,y)\in\R^2:y>0\},\qquad
 g=\frac{\mathrm dx^2+\mathrm dy^2}{y^2}.
\]
A \(C^r\) isometric immersion into \(\R^m\), \(r\ge1\), is a globally defined \(C^r\) map \(F:\mathbb H^2\to\R^m\) whose differential satisfies
\[
\langle\partial_xF,\partial_xF\rangle
=\langle\partial_yF,\partial_yF\rangle=y^{-2},
\qquad \langle\partial_xF,\partial_yF\rangle=0
\]
at every \((x,y)\in\mathbb H^2\). These equations imply rank two. No injectivity, properness, or equality between intrinsic hyperbolic distance and ambient Euclidean chord distance is required. Smooth is interpreted as \(C^\infty\).

The main question is whether such a smooth \(F:\mathbb H^2\to\R^4\) exists on the entire complete hyperbolic plane. For each specified regularity \(r\in\N_{>0}\cup\{\infty\}\), a precise version of the broader least-dimension question is to determine
\[
d_r=\inf\{m\in\N:m\ge2,\ \exists\ C^r\text{ isometric immersion }\mathbb H^2\to\R^m\},
\]
with value \(\infty\) if there is none. The phrase natural regularity assumptions does not identify further classes, such as a particular Hölder condition.

The title says embedding, whereas the actual question explicitly asks about immersion; the latter is retained. The supplied research summary reports nonexistence of \(C^2\) isometric embeddings into \(\R^3\), attributed to Efimov, and existence of \(C^1\) embeddings into \(\R^3\), attributed to Nash/Kuiper. Its embedding statements are not silently substituted for immersion statements. Its claimed smooth \(\R^4\) progress in the 2020s has no precise theorem or citation, and does not supply a verified answer to the main question or values of \(d_r\).
\subsection{Short English statement}
Does the complete hyperbolic plane have a globally defined smooth isometric immersion into \(\R^4\)? More generally, find the minimum target dimension for each specified regularity class; immersion and embedding must be distinguished.
\subsection{Sources}
\begin{itemize}
\item Local UnsolvedMath registry, permanent ID \(3900006\), exact source code \texttt{catalogue entry}, ``Embedding the hyperbolic plane.'' Proposer field: David Eppstein and cited proposers.
\item Eppstein, \emph{The Geometry Junkyard: Open Problems}, item ``Embedding the hyperbolic plane''; \url{https://ics.uci.edu/~eppstein/junkyard/open.html}. HTML extraction, recorded access 2026-07-29.
\item Supplied research summary attributes regularity-dependent embedding results to Efimov and Nash/Kuiper, and mentions unspecified 2020s progress in dimension four. No complete primary references or precise recent theorem are supplied.
\end{itemize}
\end{document}
